% Hafta 11 — Code lifting: anahtarı hiç çıkarmadan
% Derleme: py -3.12 tools/latex/derle.py docs/week-11/assets/h11-06-code-lifting.tex
\documentclass[tikz,border=8pt]{standalone}
\usepackage{cen429-cizim}
\begin{document}
\begin{tikzpicture}[node distance=8mm and 16mm]
  \node[baslik] (b) at (0,0) {Code lifting: anahtarı hiç çıkarmadan};
  \node[altbaslik, text width=150mm] at (b.south west) {WBC'nin en sinsi saldırısı: anahtar gizli kalsa bile işlev çalınır};

  \node[kotukutu=62mm, below=16mm of b.south west, anchor=north west] (sol) {\kb{Saldırı}\\[2mm]
    whitebox rutinini/tablosunu\\olduğu gibi KOPYALA\\[3mm]
    kendi programında kullan\\anahtar hiç bilinmez};
  \node[iyikutu=62mm, right=of sol] (sag) {\kb{Önlem}\\[2mm]
    dış kodlama (F, G)\\cihaz bağlama\\sunucu tarafı denetim\\[3mm]
    taşınan rutin işe yaramaz};
  \draw[draw=cizgi, line width=1.6pt] ($(sol.north east)!0.5!(sag.north west)$) -- ($(sol.south east)!0.5!(sag.south west)$);

  \node[dip, text width=160mm, below=8mm of sol.south -| sag, anchor=north] {%
    Bu yüzden WBC tek başına değil, bağlama ve sunucu denetimiyle birlikte kullanılır.};
\end{tikzpicture}
\end{document}
